%%═══════════════════════════════════════════════════════════════════
%% Lecture 07 — Stiff ODEs, Numerov Method, Boundary-Value Problems,
%%              and Relaxation
%% 77315 — Computational Physics
%% Date: 07/05/2026
%%═══════════════════════════════════════════════════════════════════

\section{Lecture 07 — Stiff ODEs, Numerov Method, Boundary-Value
  Problems, and Relaxation}\label{sec:lec07}

\begin{center}
\begin{tabular}{ll}
  \textbf{Date:}\quad 07/05/2026 & \\
\end{tabular}
\end{center}
\hrule\vspace{1.5em}

%%──────────────────────────────────────────────────────────────────
\subsection*{Overview}
Lecture~\ref{sec:lec06} ended with the stability restriction for explicit
methods applied to systems of ODEs. This lecture begins by turning that
criterion into a concrete example of a \emph{stiff} system: a system whose
fast-decaying modes force a tiny numerical step even though the physically
interesting solution evolves slowly. The lecture then introduces Numerov's
method for second-order equations of the form
$y''(x)=f(x)y(x)+g(x)$, and moves from initial-value problems to
boundary-value problems. The final part explains shooting, shooting to a
fitting point, and relaxation, with a white-dwarf structure problem as the
main physical example.

%%──────────────────────────────────────────────────────────────────
\subsection{A Stiff Two-Variable System}\label{subsec:lec07:stiff-example}

\subsubsection*{Physical Motivation}
In many physical systems several processes act on the same variables but on
very different time scales. A chemical abundance may have a very fast
transient and a slow secular evolution; a mechanical system may contain a
heavily damped mode and a slowly damped mode. Numerically, the fast mode may
control stability even after it has become physically negligible.

Consider the linear system
\begin{align}
    u' &= 998u+1998v,
    \label{eq:lec07:stiff-u}\\
    v' &= -999u-1999v,
    \label{eq:lec07:stiff-v}
\end{align}
with initial data
\begin{equation}
    u(0)=1,
    \qquad
    v(0)=0.
    \label{eq:lec07:stiff-initial-data}
\end{equation}
Introduce new variables $y$ and $z$ by
\begin{equation}
    u=2y-z,
    \qquad
    v=-y+z.
    \label{eq:lec07:stiff-change-variables}
\end{equation}
Adding the two original equations gives
\begin{equation}
    y'=u'+v'=-u-v.
    \label{eq:lec07:y-prime-start}
\end{equation}
Since $u+v=y$, this becomes
\begin{equation}
    \boxed{y'=-y}
    \label{eq:lec07:y-prime}
\end{equation}

\textit{Physical significance:} $y$ is the slow mode; it decays on the
order-one scale in $x$.

To isolate $z$, use $u+2v=z$. Therefore
\begin{equation}
    z'=u'+2v'.
    \label{eq:lec07:z-prime-start}
\end{equation}
Substituting Eqs.~\eqref{eq:lec07:stiff-u}--\eqref{eq:lec07:stiff-v}
and then Eq.~\eqref{eq:lec07:stiff-change-variables} gives
\begin{align}
    z'
    &= (998u+1998v)+2(-999u-1999v)
    \label{eq:lec07:z-prime-algebra-one}\\
    &= -1000u-2000v
    = -1000(u+2v)
    = -1000z.
    \label{eq:lec07:z-prime-algebra-two}
\end{align}
Thus
\begin{equation}
    \boxed{z'=-1000z}
    \label{eq:lec07:z-prime}
\end{equation}

\textit{Physical significance:} $z$ is the fast transient; it decays on the
scale $10^{-3}$ in $x$.

The initial data imply
\begin{equation}
    y(0)=u(0)+v(0)=1,
    \qquad
    z(0)=u(0)+2v(0)=1.
    \label{eq:lec07:yz-initial-data}
\end{equation}
Therefore
\begin{equation}
    y(x)=e^{-x},
    \qquad
    z(x)=e^{-1000x}.
    \label{eq:lec07:yz-solutions}
\end{equation}
Returning to $u$ and $v$,
\begin{align}
    \boxed{u(x)=2e^{-x}-e^{-1000x}},
    \label{eq:lec07:u-solution}\\
    \boxed{v(x)=-e^{-x}+e^{-1000x}}.
    \label{eq:lec07:v-solution}
\end{align}

\textit{Physical significance:} after a very short initial layer, the
solution is dominated by $e^{-x}$, but an explicit integrator still has to
respect the stability constraint imposed by the $e^{-1000x}$ mode.

For an explicit method such as Euler, Lecture~\ref{sec:lec06} showed that a
decay mode $e^{-\lambda x}$ imposes roughly
\begin{equation}
    h<\frac{2}{\lambda}.
    \label{eq:lec07:euler-mode-stability}
\end{equation}
Here the largest decay rate is $\lambda=1000$, so
\begin{equation}
    \boxed{h<2\times 10^{-3}}
    \label{eq:lec07:stiff-step-restriction}
\end{equation}

\textit{Physical significance:} the numerical method must take thousands of
tiny steps per unit interval even though the surviving physical mode varies
on a scale of order $1$. This mismatch is the hallmark of stiffness.

Implicit or semi-implicit methods are designed precisely for this situation:
they damp the fast stable modes without forcing the step size to resolve
their decay in detail.

%%──────────────────────────────────────────────────────────────────
\subsection{Numerov's Method for Second-Order Equations}
\label{subsec:lec07:numerov}

\subsubsection*{Physical Motivation}
Many physics problems reduce directly to a second-order equation rather than
to an arbitrary first-order system. If the equation is linear in $y$, then
we can exploit the special structure to obtain a very accurate three-point
method with little extra work. This is the purpose of Numerov's method.

Consider
\begin{equation}
    y''(x)=f(x)y(x)+g(x),
    \label{eq:lec07:numerov-equation}
\end{equation}
where $f$ and $g$ are known functions. Two common physical examples are:
\begin{align}
    \psi''(x) &= -k^2(x)\psi(x),
    \label{eq:lec07:wave-example}\\
    \frac{1}{r^2}\frac{d}{dr}
    \left(r^2\frac{d\phi}{dr}\right)
    &= -4\pi\rho(r).
    \label{eq:lec07:poisson-radial}
\end{align}
In the radial Poisson problem, defining $u(r)=r\phi(r)$ converts the radial
operator to a second derivative:
\begin{equation}
    u''(r)=-4\pi r\rho(r).
    \label{eq:lec07:poisson-u}
\end{equation}
This is a special case of Eq.~\eqref{eq:lec07:numerov-equation}.

\subsubsection*{Physical Motivation}
The central finite-difference formula for $y''$ is second order. Numerov's
method improves it by using the differential equation itself to cancel the
leading fourth-derivative error. The result is a symmetric stencil involving
$x-h$, $x$, and $x+h$.

Expand $y(x+h)$ and $y(x-h)$ around $x$:
\begin{align}
    y(x+h)
    &=y+h y'+\frac{h^2}{2}y''
      +\frac{h^3}{6}y'''
      +\frac{h^4}{24}y^{(4)}
      +\frac{h^5}{120}y^{(5)}
      +\bigO{h^6},
    \label{eq:lec07:taylor-plus}\\
    y(x-h)
    &=y-h y'+\frac{h^2}{2}y''
      -\frac{h^3}{6}y'''
      +\frac{h^4}{24}y^{(4)}
      -\frac{h^5}{120}y^{(5)}
      +\bigO{h^6}.
    \label{eq:lec07:taylor-minus}
\end{align}
Adding them cancels the odd derivatives:
\begin{equation}
    y(x+h)+y(x-h)
    =2y+h^2y''+\frac{h^4}{12}y^{(4)}+\bigO{h^6}.
    \label{eq:lec07:taylor-sum}
\end{equation}
Now apply the same central sum to $y''$:
\begin{equation}
    y''(x+h)+y''(x-h)
    =2y''+h^2y^{(4)}+\bigO{h^4}.
    \label{eq:lec07:second-derivative-sum}
\end{equation}
Multiplying Eq.~\eqref{eq:lec07:second-derivative-sum} by $h^2/12$ and
subtracting it from Eq.~\eqref{eq:lec07:taylor-sum} eliminates
$y^{(4)}$:
\begin{equation}
    y(x+h)+y(x-h)
    -\frac{h^2}{12}\left[y''(x+h)+y''(x-h)\right]
    =2y+\frac{5h^2}{6}y''+\bigO{h^6}.
    \label{eq:lec07:numerov-before-substitution}
\end{equation}
Let
\begin{equation}
    x_i=x_0+ih,
    \qquad
    y_i=y(x_i),
    \qquad
    f_i=f(x_i),
    \qquad
    g_i=g(x_i).
    \label{eq:lec07:grid-definitions}
\end{equation}
Substitute $y_i''=f_i y_i+g_i$ into
Eq.~\eqref{eq:lec07:numerov-before-substitution}. After collecting the
terms with $y_{i+1}$, $y_i$, and $y_{i-1}$, one obtains
\begin{equation}
    \boxed{
    \left(1-\frac{h^2}{12}f_{i+1}\right)y_{i+1}
    -2\left(1+\frac{5h^2}{12}f_i\right)y_i
    +\left(1-\frac{h^2}{12}f_{i-1}\right)y_{i-1}
    =
    \frac{h^2}{12}
    \left(g_{i+1}+10g_i+g_{i-1}\right)}
    \label{eq:lec07:numerov-formula}
\end{equation}
with local error $\bigO{h^6}$.

\textit{Physical significance:} Numerov's method advances a linear
second-order equation with high accuracy while requiring only one new
evaluation of the known coefficients at the next grid point.

Because Eq.~\eqref{eq:lec07:numerov-formula} is symmetric, it naturally uses
equally spaced points. The step size can still be changed in factors of two:
to halve $h$, compute an intermediate point and continue on the finer grid;
to double $h$, skip every other point in the next stencil. Like all
multi-point methods, Numerov also needs a startup procedure, because the
values at two neighboring points must be known before the recurrence can
begin.

%%──────────────────────────────────────────────────────────────────
\subsection{Boundary-Value Problems and Shooting}
\label{subsec:lec07:boundary-value}

\subsubsection*{Physical Motivation}
An initial-value problem gives all needed data at one endpoint and asks us
to march forward. Many physical problems instead specify conditions at two
different places: a wavefunction must be regular at the origin and decay at
infinity, a string may be fixed at both ends, and a stellar model must be
regular at the center and have vanishing pressure at the surface. These are
boundary-value problems.

For a second-order equation written as a first-order system,
\begin{align}
    \frac{dy}{dx} &= F(x,y,z),
    \label{eq:lec07:bvp-system-y}\\
    \frac{dz}{dx} &= G(x,y,z),
    \label{eq:lec07:bvp-system-z}
\end{align}
an initial-value problem might specify both $y(x_1)$ and $z(x_1)$. A
boundary-value problem may instead specify
\begin{equation}
    y(x_1)=a,
    \qquad
    y(x_2)=b,
    \label{eq:lec07:bvp-simple-boundaries}
\end{equation}
leaving the initial derivative-like variable $z(x_1)$ unknown.

The shooting method treats the missing initial datum as a parameter. Choose
a trial value
\begin{equation}
    z(x_1)=s,
    \label{eq:lec07:shooting-trial-slope}
\end{equation}
integrate the initial-value problem to $x_2$, and define the mismatch
\begin{equation}
    \Delta(s)=y(x_2;s)-b.
    \label{eq:lec07:shooting-mismatch}
\end{equation}
The correct initial slope satisfies
\begin{equation}
    \boxed{\Delta(s)=0}
    \label{eq:lec07:shooting-root}
\end{equation}

\textit{Physical significance:} shooting converts a boundary-value problem
into a root-finding problem for the unknown initial data.

For an $N$-dimensional first-order system
\begin{equation}
    \frac{dy_i}{dx}
    =f_i(x,y_1,\ldots,y_N),
    \qquad i=1,\ldots,N,
    \label{eq:lec07:general-first-order-system}
\end{equation}
suppose $N_1$ boundary conditions are imposed at $x_1$ and $N_2$ are imposed
at $x_2$, with
\begin{equation}
    N=N_1+N_2.
    \label{eq:lec07:condition-count}
\end{equation}
The boundary conditions need not be simple values. They may be relations
\begin{align}
    B_j^{(1)}(x_1,y_1,\ldots,y_N)&=0,
    \qquad j=1,\ldots,N_1,
    \label{eq:lec07:left-boundary-relations}\\
    B_j^{(2)}(x_2,y_1,\ldots,y_N)&=0,
    \qquad j=1,\ldots,N_2.
    \label{eq:lec07:right-boundary-relations}
\end{align}
Shooting guesses the $N_2$ missing degrees of freedom at $x_1$, integrates
to $x_2$, and solves the $N_2$ mismatch equations there.

\textit{Physical significance:} the number of guessed initial quantities
must match the number of missing right-boundary constraints; otherwise the
boundary-value problem is either underdetermined or overdetermined.

%%──────────────────────────────────────────────────────────────────
\subsection{Shooting to a Fitting Point}\label{subsec:lec07:fitting-point}

\subsubsection*{Physical Motivation}
Sometimes integration from the left endpoint all the way to the right is
numerically unstable or impossible because of singular behavior. In such
cases it is better to integrate from both ends toward a common point and
demand that the two partial solutions meet there.

Let $x_m$ be an interior fitting point. Integrate from $x_1$ to $x_m$ using
unknown parameters $\mathbf{s}_L$, and integrate from $x_2$ back to $x_m$
using unknown parameters $\mathbf{s}_R$. The matching condition is
\begin{equation}
    \boxed{
    \mathbf{y}_L(x_m;\mathbf{s}_L)
    -\mathbf{y}_R(x_m;\mathbf{s}_R)=\mathbf{0}}
    \label{eq:lec07:fitting-condition}
\end{equation}

\textit{Physical significance:} shooting to a fitting point avoids forcing a
single integration through a region where the numerical solution would
become badly conditioned.

This method is especially useful when each boundary selects a locally stable
branch of the solution. Each side is integrated in the direction in which it
behaves well, and the root-finding problem adjusts the free parameters until
the two branches join smoothly.

%%──────────────────────────────────────────────────────────────────
\subsection{White-Dwarf Structure as a Boundary-Value Problem}
\label{subsec:lec07:white-dwarf}

\subsubsection*{Physical Motivation}
A white dwarf is supported by degeneracy pressure. In the simplified model
used here, the pressure depends only on density and not on temperature. The
star is therefore a clean example of a boundary-value problem: gravity pulls
inward, the pressure gradient pushes outward, the center must be regular,
and the pressure must vanish at the surface.

Use the enclosed mass $m$ as the independent variable. Hydrostatic
equilibrium may be written as
\begin{equation}
    \frac{dP}{dm}
    =-\frac{Gm}{4\pi r^4},
    \label{eq:lec07:hydrostatic-mass}
\end{equation}
where $P$ is pressure, $G$ is Newton's constant, and $r(m)$ is the radius
enclosing mass $m$. The equation of state is approximated by a polytrope,
\begin{equation}
    \boxed{P=K\rho^\gamma}
    \label{eq:lec07:polytropic-eos}
\end{equation}

\textit{Physical significance:} the exponent $\gamma$ measures how strongly
the pressure rises when the star is compressed.

For nonrelativistic degenerate matter,
\begin{equation}
    \gamma=\frac{5}{3},
    \label{eq:lec07:gamma-nonrel}
\end{equation}
while in the relativistic limit,
\begin{equation}
    \gamma=\frac{4}{3}.
    \label{eq:lec07:gamma-rel}
\end{equation}
The lecture used an order-of-magnitude estimate to understand the
mass-radius relation. Approximate the central pressure scale by
$P/M$ in Eq.~\eqref{eq:lec07:hydrostatic-mass}, and estimate the density by
\begin{equation}
    \rho\sim \frac{M}{\frac{4\pi}{3}R^3},
    \label{eq:lec07:mean-density}
\end{equation}
where $M$ and $R$ are the total mass and outer radius. Then
\begin{equation}
    \frac{K\rho^\gamma}{M}
    \sim \frac{GM}{4\pi R^4}.
    \label{eq:lec07:scaling-balance-start}
\end{equation}
Substituting Eq.~\eqref{eq:lec07:mean-density} gives the scaling relation
\begin{equation}
    K\left(\frac{M}{R^3}\right)^\gamma
    \sim \frac{GM^2}{R^4}.
    \label{eq:lec07:scaling-balance}
\end{equation}
Rearranging,
\begin{equation}
    R^{4-3\gamma}\sim
    \frac{G}{K}M^{2-\gamma}.
    \label{eq:lec07:mass-radius-general}
\end{equation}
For $\gamma=5/3$,
\begin{equation}
    \boxed{R\propto M^{-1/3}}
    \label{eq:lec07:mass-radius-nonrel}
\end{equation}

\textit{Physical significance:} a more massive nonrelativistic white dwarf
is smaller, because stronger gravity requires a higher density and therefore
a smaller radius.

For $\gamma=4/3$, the radius drops out of the scaling relation:
\begin{equation}
    M^{2-\gamma}=M^{2/3}\sim \text{constant}.
    \label{eq:lec07:mass-limit-scaling}
\end{equation}
Thus the model predicts a special mass scale rather than a free
mass-radius curve:
\begin{equation}
    \boxed{M\sim M_{\mathrm{Ch}}}
    \label{eq:lec07:chandrasekhar-scale}
\end{equation}

\textit{Physical significance:} in the relativistic degenerate limit,
increasing density no longer raises pressure fast enough to balance gravity
for arbitrary mass. This is the origin of the Chandrasekhar-limit behavior.

The estimate is only qualitative, because it replaces a differential
equation by a global scaling argument. The actual numerical task is to solve
the discretized hydrostatic equilibrium equations.

%%──────────────────────────────────────────────────────────────────
\subsection{Relaxation Discretization for the White Dwarf}
\label{subsec:lec07:relaxation}

\subsubsection*{Physical Motivation}
Shooting changes a few initial parameters and repeatedly integrates across
the domain. Relaxation takes a more global view: guess the entire solution
on a grid and then adjust all grid values until every discrete equation and
boundary condition is satisfied at once.

Divide the star into $N$ mass shells. Let
\begin{equation}
    m_i=\frac{i}{N}M,
    \qquad i=1,\ldots,N,
    \label{eq:lec07:mass-grid}
\end{equation}
and use the shell radii $r_i$ as the unknowns. The center condition is
\begin{equation}
    r_0=0,
    \qquad
    m_0=0.
    \label{eq:lec07:center-conditions}
\end{equation}
The pressure in shell $i$ is
\begin{equation}
    P_i=K\rho_i^\gamma,
    \label{eq:lec07:shell-pressure}
\end{equation}
where the shell density is mass divided by shell volume:
\begin{equation}
    \rho_i=
    \frac{m_i-m_{i-1}}
    {\frac{4\pi}{3}\left(r_i^3-r_{i-1}^3\right)}.
    \label{eq:lec07:shell-density}
\end{equation}
The finite-difference pressure gradient around shell interface $i$ is
\begin{equation}
    \frac{\Delta P_i}{\Delta m_i}
    =
    \frac{P_{i+1}-P_i}
    {\frac{1}{2}(m_{i+1}+m_i)-\frac{1}{2}(m_i+m_{i-1})}.
    \label{eq:lec07:pressure-gradient-discrete}
\end{equation}
For the uniform mass grid this denominator is simply $M/N$.

Define the residual
\begin{equation}
    F_i(\mathbf{r})
    =
    \frac{\Delta P_i}{\Delta m_i}
    \left(\frac{4\pi r_i^4}{Gm_i}\right)
    +1,
    \qquad i=1,\ldots,N.
    \label{eq:lec07:relaxation-residual}
\end{equation}
The discrete hydrostatic equilibrium problem is
\begin{equation}
    \boxed{
    F_i(\mathbf{r})=0,
    \qquad i=1,\ldots,N}
    \label{eq:lec07:relaxation-system}
\end{equation}

\textit{Physical significance:} each residual measures how far one mass
shell is from balancing the pressure gradient against gravity.

At the outer boundary the pressure outside the star is zero. This is encoded
by
\begin{equation}
    P_{N+1}=0,
    \qquad
    m_{N+1}=m_N=M.
    \label{eq:lec07:surface-conditions}
\end{equation}
Together with the center conditions, the $N$ equations determine the
$N$ unknown radii $r_1,\ldots,r_N$.

Newton's method for the vector residual is
\begin{equation}
    \sum_{j=1}^{N}
    \frac{\partial F_i}{\partial r_j}\Delta r_j
    =-F_i(\mathbf{r}),
    \qquad i=1,\ldots,N.
    \label{eq:lec07:newton-linear-system}
\end{equation}
After solving this linear system, update
\begin{equation}
    r_i^{\mathrm{new}}=r_i+\Delta r_i.
    \label{eq:lec07:newton-radius-update}
\end{equation}
The derivatives $\partial F_i/\partial r_j$ should be evaluated
analytically, because the residuals depend only on neighboring shell radii
and the resulting Jacobian is structured.

\textit{Physical significance:} relaxation solves for the whole stellar
profile at once; the price is that each iteration requires a linear solve
with the Jacobian matrix.

The initial guess requested in the lecture is a linear radius profile,
\begin{equation}
    r_i^{(0)}=\frac{i}{N}R_{\mathrm{guess}},
    \label{eq:lec07:linear-radius-guess}
\end{equation}
with $N=100$ for the main calculations. Because the first guess can be far
from the solution, the Newton correction must be damped if it would make the
radius grid nonmonotonic. Use
\begin{equation}
    r_i^{\mathrm{new}}=r_i+\alpha\Delta r_i,
    \qquad
    0<\alpha\leq 1,
    \label{eq:lec07:damped-newton-update}
\end{equation}
and choose $\alpha$ small enough that
\begin{equation}
    r_{i+1}^{\mathrm{new}}>r_i^{\mathrm{new}}
    \qquad
    \text{for all } i.
    \label{eq:lec07:monotone-radius-condition}
\end{equation}

\textit{Physical significance:} shell crossing would imply negative or
infinite densities, so damping the Newton step preserves the physical
ordering of the mass shells.

The lecture also emphasizes a practical lesson: a black-box nonlinear solver
from a library may fail when the initial guess is too far from the physical
solution. A custom relaxation algorithm can include problem-specific damping
and monotonicity checks, which makes it more robust for this stellar
structure problem.

%%──────────────────────────────────────────────────────────────────
\subsection{Computational Tasks from the Lecture}
\label{subsec:lec07:computational-tasks}

\subsubsection*{Physical Motivation}
The homework tasks are designed to connect the numerical algorithms to a
real physical model. The point is not only to obtain numbers, but also to
test whether the numerical method reproduces the expected scaling laws and
where the simplified physics begins to fail.

The requested solver should:
\begin{enumerate}
    \item implement a general Newton solver for $N$ equations in $N$
    unknowns, using both a residual function $\mathbf{F}(\mathbf{x})$ and an
    analytic Jacobian $\partial F_i/\partial x_j$;
    \item test this solver on a problem with a known analytic solution;
    \item implement the white-dwarf residuals
    Eq.~\eqref{eq:lec07:relaxation-residual} and their analytic Jacobian;
    \item solve the nonrelativistic case $\gamma=5/3$ for
    $M=1M_\odot$, $R_{\mathrm{guess}}=1R_\odot$, and $N=100$;
    \item plot and tabulate $\rho(r)$ in CGS units;
    \item solve several masses and check whether $RM^{1/3}$ is
    approximately constant;
    \item solve the nearly relativistic case $\gamma=1.334$ for masses near
    $1.4M_\odot$ and inspect how the outer radius changes;
    \item compare with a SciPy nonlinear solver, noting that it may require
    an initial radius guess closer to the final solution.
\end{enumerate}

The central diagnostic for the nonrelativistic sequence is
\begin{equation}
    \boxed{
    RM^{1/3}\approx \text{constant}}
    \label{eq:lec07:nonrel-diagnostic}
\end{equation}

\textit{Physical significance:} this numerical check tests whether the full
discrete hydrostatic model reproduces the scaling argument
$R\propto M^{-1/3}$.

For the nearly relativistic sequence, small changes in mass can produce
large changes in radius:
\begin{equation}
    \gamma=1.334,
    \qquad
    M/M_\odot=1.40,1.41,\ldots,1.46.
    \label{eq:lec07:near-relativistic-sequence}
\end{equation}

\textit{Physical significance:} the sensitivity near $\gamma=4/3$ reflects
the approach to the limiting-mass behavior discussed in
Eq.~\eqref{eq:lec07:chandrasekhar-scale}.

%%──────────────────────────────────────────────────────────────────
\subsection*{Recommended Reading}
\begin{itemize}
    \item \textit{Computational Physics}, Mark Newman --- Ch.\ 8.6--8.8
    (boundary-value problems and relaxation ideas).
    \item \textit{Numerical Recipes}, W.H. Press et al.\ --- Ch.\ 17.4 and
    Ch.\ 18.1--18.2 (stiff equations, shooting, relaxation, and boundary
    value problems).
    \item Tao Pang, \textit{An Introduction to Computational Physics} ---
    Ch.\ 3.5--3.6 (Numerov's method and boundary-value problems).
    \item S. Chandrasekhar, \textit{An Introduction to the Study of Stellar
    Structure} --- classic background on polytropes and white-dwarf
    mass-radius relations.
\end{itemize}
